\section{Introduction}
\label{sec:intro}

The United States Constitution establishes congressional representation through
two distinct clauses. Article~I, \S2 mandates that seats be
\emph{apportioned among the several States} according to population, a
requirement implemented since 1941 by the Huntington-Hill method
\citep{balinski2001fair}. Article~I, \S4 provides that ``the Times, Places and
Manner of holding Elections \ldots\ shall be prescribed in each State by the
Legislature thereof,'' but that ``Congress may at any time by make or alter
such Regulations.'' In practice, these clauses are treated as governing
separate questions: \S2 determines \emph{how many} seats each state receives;
\S4 (via state or congressional action) determines \emph{where} the district
boundaries lie.

We argue this separation is artificial and consequential. The seat count $n$
assigned to a state is not merely a quota --- it encodes a geometric structure.
Specifically, the prime factorization $n = p_1^{a_1} \times p_2^{a_2} \times
\cdots$ determines a unique hierarchical partition of any geographic region:
split into $p_1$ sub-regions, then split each sub-region into $p_2$ parts, and
so on. When applied to a state's census-tract graph using minimum-cut graph
partitioning, this hierarchy produces a congressional map with two inputs only:
\begin{enumerate}
  \item The prime factorization of $n$ --- pure arithmetic, determined by the
        apportionment process.
  \item The census-tract adjacency graph with TIGER boundary weights --- pure
        geography, derived from federal public records.
\end{enumerate}
No political data enters. No human boundary choices are made. The resulting
\emph{ApportionRegions} (AR) algorithm is a direct geographic
extension of Huntington-Hill: the same constitutional process that decides how
many seats a state receives also decides, without further human input, where
those seats are located.

\paragraph{The reuse property.}
The most practically significant feature of AR is that geographic splits are
computed once and reused. The canonical 2-way minimum-cut partition of a state
(``the bisection'') is computed once and applies to every even seat count.
The canonical 3-way partition applies to every seat count divisible by~3.
When a state gains or loses a seat after reapportionment, only the levels of
the factorization tree affected by the changed prime structure require
recomputation. Districts have \emph{ancestors} in the tree --- a district
in the $n$-seat map can be traced to the leaf of the factorization tree that
contains its census tracts, and that leaf's ancestry is preserved across all
seat counts sharing the same prime structure at that level.

\paragraph{The discontinuity problem.}
AR also makes visible a previously unrecognized structural phenomenon:
consecutive seat counts can have disjoint prime factorizations, causing
complete structural breaks in the map. California currently has 52 seats
($52 = 2^2 \times 13$); the preceding apportionment gave it 53 seats
($53$ is prime). Going from 52 to 53 requires discarding the entire
2-level binary structure and replacing it with a single 53-way cut of the
whole state. Going from 51 ($= 3 \times 17$) to 52 ($= 4 \times 13$)
replaces a ternary top level with a binary one. These discontinuities are
\emph{mathematically forced}, not arbitrary: they reflect the prime
structure of the integers, not political choices. We characterize when such
breaks occur and propose a continuity-preserving variant (AR-smooth) that
minimises structural change across reapportionments.

\paragraph{Contributions.}
\begin{enumerate}
  \item \textbf{Algorithm}: The AR algorithm --- a deterministic,
        prime-factorization-guided redistricting procedure with formal
        reuse semantics and a two-input specification (prime factorization
        + census geography).
  \item \textbf{Constitutional argument}: A formal claim that AR maps
        arise from Article~I~\S2 (apportionment) rather than \S4 (Manner),
        shifting the constitutional footing and removing the ``separate
        districting question.''
  \item \textbf{Reuse theorem}: A proof that prime-$p$ splits are
        canonical and reusable, and that the factorization tree is unique
        given the minimum-cut at each level.
  \item \textbf{Empirical evaluation}: Full AR maps for all 50 states at
        their 2020 apportionment seat counts, comparison to minimum-edge-cut
        (MEC) maps from \citet{dellaLibera2026mec}, and characterisation of
        the 51$\to$52 California discontinuity.
  \item \textbf{AR-smooth}: A reapportionment-stable variant that preserves
        as much of the existing factorization tree as possible when $n$
        changes by~$\pm 1$.
\end{enumerate}
